\documentclass[11pt]{article}
\usepackage[a4paper,margin=19mm]{geometry}
\usepackage{fontspec}
\setmainfont{lmroman10-regular.otf}
\usepackage{amsmath,amssymb,amsthm,mathtools,amscd,xcolor,graphicx,hyperref}
\hypersetup{hidelinks}
\numberwithin{equation}{section}
\newtheorem{theorem}{Theorem}[section]
\newtheorem{lemma}[theorem]{Lemma}
\newtheorem{proposition}[theorem]{Proposition}
\newtheorem{corollary}[theorem]{Corollary}
\newtheorem{claim}[theorem]{Claim}
\newtheorem{fact}[theorem]{Fact}
\newtheorem{conjecture}[theorem]{Conjecture}
\theoremstyle{definition}
\newtheorem{definition}[theorem]{Definition}
\newtheorem{example}[theorem]{Example}
\newtheorem{remark}[theorem]{Remark}
\newtheorem{question}[theorem]{Question}
\theoremstyle{plain}
\newcommand{\C}{\mathbb C}
\newcommand{\R}{\mathbb R}
\newcommand{\Q}{\mathbb Q}
\newcommand{\Z}{\mathbb Z}
\newcommand{\N}{\mathbb N}
\newcommand{\eee}{\mathrm e}
\newcommand{\cclass}[1]{\textsf{#1}}
\newcommand{\cproblem}[1]{\textsc{#1}}
\newcommand{\ie}{i.\,e.}
\newcommand{\Ie}{I.\,e.}
\newcommand{\eg}{e.\,g.}
\newcommand{\Eg}{E.\,g.}
\newcommand{\phd}{Ph.\,D.}
\newcommand{\msc}{M.\,S.}
\newcommand{\bsc}{B.\,S.}
\newcommand{\expref}[2]{\hyperref[#2]{#1~\ref{#2}}}
\newcommand{\expeqref}[2]{\hyperref[#2]{#1~\eqref{#2}}}
\newcommand{\exprefsilent}[2]{\hyperref[#2]{#1}}
\providecommand{\etal}{et~al.}
\usepackage{eprint}
\usepackage{packages,aumacros}
\title{Strong XOR direct sum for expected randomized query complexity of partial functions}
\author{Joshua Brody, Jae Tak Kim, Peem Lerdputtipongporn, Hariharan Srinivasulu}
\date{doi:10.4086/toc.2023.v019a011; CC BY 3.0}
\begin{document}
\maketitle
\noindent\textit{Editorial note.} Complete original author text, definitions, proof and bibliography retained. Expected versus worst-case query costs and original error/abort parameters preserved. Standalone wrapper only.\par\medskip
\noindent\textit{Allowed standard prerequisites (editorial).}\begin{itemize}
\item Randomized decision trees and elementary product probability. All original expected-cost, distributional and aborting-query definitions and local lemmas retained.. No collector proof or strengthening from expected to worst-case query cost.
\end{itemize}\medskip
\section{Introduction}

We show that $\xor$ admits a strong direct sum theorem for randomized
query complexity.  Generally, the \emph{direct sum problem} asks how
the cost of computing a
partial %
function $g$ scales with the number $k$ of
instances of the function that we need to compute
simultaneously (in parallel) %
.
%
%
This is a
foundational computational problem that has received considerable
attention~\cite{ImpagliazzoRW94,BenAsherN95,NisanRS99, Shaltiel03,
  JainKS10, BenDavidK18,Drucker12, BlaisB19, BenDavidB20a, BenDavidB20b,
  BenDavidGKW20}, including recent %
a recent paper by %
Blais and
Brody~\cite{BlaisB19}, which showed that %
\emph{expected} %
%
query complexity obeys a direct sum theorem in a strong
%
sense---computing  %
$k$ copies of a
partial %
function $g$ with overall error
$\eps$ requires $k$ times the cost of computing $g$ on one input with
very low %
%
$(\eps/k)$ error.  %
This matches the naive success
amplification algorithm which runs an $\frac{\eps}{k}$-error algorithm for
$g$ once on each of $k$ inputs and applies a union bound to get an
overall error guarantee of $\eps$.

What happens if we %
do not %
need to compute $g$ on all instances, but
only on a \emph{function} $f\circ g$ of those instances?  Clearly the
same success amplification trick (compute $g$ on each input with low
error, then apply $f$ to the answers) works for computing $f\circ g$;
however, in principle, computing $f\circ g$ can be easier than
computing each instance of $g$ individually.  When a function
$f\circ g$ requires success amplification for all $g$, we say that $f$
\emph{admits a strong direct sum theorem}.  Our main result shows that
$\xor$ admits a strong direct sum theorem.


%
\subsection{Query complexity}  %
%
A \emph{query algorithm},  %
also known as a \emph{decision tree},  %
computing $f$, is an algorithm $\cA$ that takes %
an input $x$ to $f$, examines (or \emph{queries}) bits of $x$, and
outputs an answer for $f(x)$.  A \emph{leaf} of $\cA$ is a bit string
$q \in \bit^*$ representing the answers to the queries made by $\cA$
on input $x$.  Let $\leaf(\cA,x)$ denote the leaf of $\cA$ reached on input
$x$.  Naturally, our general goal is to minimize the 
length of $q$, %
\ie, minimize the number of queries needed to compute $f$.

A randomized algorithm $\cA$ \emph{computes} a function
$f : \bitsn \to \bit$ \emph{with error $\epsilon \ge 0$} if for every
input $x \in \bitsn$, the algorithm outputs the value $f(x)$ with
probability at least $1-\epsilon$.  The \emph{query cost} of $\cA$ is
the maximum number of bits of $x$ that it queries, with the
maximum taken over both the choice of input $x$ and the internal
randomness of $\cA$.  The \emph{$\epsilon$-error %
  randomized query complexity} of $f$ (also known as the
\emph{randomized decision tree complexity} of $f$) is the minimum
query cost of an algorithm $\cA$ that computes $f$ with error at
most $\epsilon$. We denote this complexity by $\rdt_\epsilon(f)$, and
we write $\rdt(f) := \rdt_{\frac13}(f)$ to denote the $\frac13$-error
randomized query complexity of $f$.

Another natural measure for the query cost of a randomized algorithm
$\cA$ is the \emph{expected} number of coordinates of an input $x$
that it queries.   %
%
%
%
Taking the maximum expected number of coordinates
queried by $\cA$ over all inputs yields the %
\emph{expected query cost} %
of $\cA$. The minimum %
expected  %
query cost of an algorithm $\cA$
that computes a function $f$ with error at most $\epsilon$ is the
%
\emph{$\epsilon$-error expected query complexity} %
of $f$, which we
denote by $\rdta_{\epsilon}(f)$.   %
%
%
%
We again write $\rdta(f) := \rdta_{\frac 13}(f)$. Note that
$\rdta_0(f)$ corresponds to the standard notion of
\emph{zero-error randomized query complexity} of $f$.
%
%
%
%
%

%
\subsection{Our results}

Our main result is a strong direct sum theorem for $\xor$.


\begin{theorem}\label{thm:sds-xor}
  For every
  partial %
  function $g:\bitsn \rightarrow \bit$ and all $\eps>0$, we
  have $\rdta_\eps(\xorg) = \Omega(k\cdot \rdta_{\eps/k}(g)).$
\end{theorem}


%
%
%
%
%
%
%
%
%
%
%
%
%
%
 This answers Conjecture~1 of Blais and
Brody~\cite{BlaisB19} in the affirmative.  We prove

\expref{Theorem}{thm:sds-xor} by proving an analogous result in
distributional query complexity.  %
%
%
We also allow our algorithms to \emph{abort} with 
%
%
%
%
%
a given  %
probability.
Let $\mu$ be a distribution on valid inputs for $f$. %
Let $\ddt_{\delta,\eps}^\mu(f)$ denote the minimal query cost of a
deterministic query algorithm that aborts with probability at most
$\delta$ and errs with probability at most $\eps$, where the
probability is taken over inputs $X \sim \mu$.  %
%
Similarly, let
$\rdt_{\delta,\eps}(f)$ denote the minimal query cost of a randomized
algorithm that computes $f$ with abort probability at most $\delta$
and error probability at most $\eps$ %
for each
valid %
input. %
%
(Here the  %
probabilities are taken over the internal randomness of 
the algorithm.)

Our main technical result is the following strong direct sum result
for $\xorg$ for distributional algorithms.
\begin{lemma}[Main Technical Lemma, informally stated.]\label{lemma:main}
  For every
  partial %
  function $g:\bitsn \rightarrow \bit$, every distribution
  $\mu$
  on the set of valid inputs  %
  and every
  %
  sufficiently small  %
  $\delta,\eps>0$, we have
  $$\ddt_{\delta,\eps}^{\mu^k}(\xorg) =
  \Omega(k\ddt_{\delta',\eps'}^\mu(g))\ ,$$ for $\delta' = \Theta(1)$
  and $\eps' = \Theta(\eps/k)$.
  %
\end{lemma}

In~\cite{BlaisB19}, Blais and Brody also gave a total function
$g:\bitsn\rightarrow \bit$ whose $\eps$-error
%
expected  %
query complexity
satisfies $\rdta_\eps(g) = \Omega(\rdt(g)\cdot \log \frac{1}{\eps})$.
We use our strong XOR Lemma together with this function to show the following.

\begin{corollary}\label{cor:application}
  There exists a total function $g:\bitsn \rightarrow \bit$ such that
  $$\rdt_\eps(\xorg) = \Omega(k\log(k)\cdot \rdt_\eps(g))\ .$$
\end{corollary}
\begin{proof}
  Let $g: \bitsn \rightarrow \bit$ be a function guaranteed
  by~\cite{BlaisB19}.  Then, we have
  $$\rdt(\xorg) \geq \rdta(\xorg) \geq \Omega(k\cdot \rdta_{1/3k}(g)) \geq
  \Omega(k\cdot \rdt(g)\cdot \log(3k)) = \Omega(k\log(k)\cdot \rdt(g))\ ,$$
  where the second inequality is by \expref{Theorem}{thm:sds-xor} and the
  third inequality is from the query complexity guarantee of $g$.
\end{proof}


This answers Open Question 1 from %
a recent paper by  %
Ben-David et al.~\cite{BenDavidGKW20}.

%
\subsection{Previous and related work}

Jain et al.~\cite{JainKS10} gave direct sum theorems for deterministic
and randomized query complexity.  %
%
%
%
%
In particular, Jain et al. show
$\rdt_\eps(f^k) \geq \delta \cdot k \cdot \rdt_{\eps + \delta}(f)$.
While their direct sum result holds for randomized query complexity,
the lower bound is in terms of the query complexity of computing $f$
with an \emph{increased error} of $\eps+\delta$.  This weakens the right-hand side of their inequality.
%
%
%
%
%
%
Shaltiel~\cite{Shaltiel03} gave a function $f$ such %
that $\ddt_{0,\eps}^{\mu^k}(f^k) \ll k\ddt_{0,\eps}^\mu(f)$, thus %
showing that a similar direct sum theorem fails to hold for %
distributional complexity. %

%
%
%
%
%
%
%
%
Drucker~\cite{Drucker12} gave a
\emph{direct product theorem} for randomized query complexity, showing
that any algorithm computing $g^k$ using $\alpha k\rdt(g)$ queries for
a constant $\alpha<1$ has success probability exponentially small in
$k$.  %
Drucker also gave the following XOR Lemma, showing that any
algorithm for $\xorg$ that makes $\ll kR(g)$ queries has
success probability exponentially close to $1/2$
%
%
%
\cite[Theorem 1.3]{Drucker12}.
%
\begin{theorem}[Drucker]
  %
  %
  %
Suppose any
%
randomized %
$T$-query algorithm has success
probability  %
$\leq 1-\eps'$ in computing
%
the Boolean function %
$g$ on input $x \sim \mu$
for some input distribution $\mu$.
%
Then, for all $0<\alpha<1$, any randomized  
algorithm making $\alpha \eps' Tk$ queries to compute $\xorg$
on input distribution $\mu^k$   %
($k$ inputs drawn independently from $\mu$)  %
has success probability at most
  $\frac{1}{2}\left(1 + [1-2\eps' +
    6\alpha\ln(2/\alpha)\eps']^k\right)$.
\end{theorem}

Drucker's XOR Lemma applies to %
%
%
%
randomized  %
query complexity 
$\rdt(\xorg)$, while ours applies to %
expected %
randomized  %
query complexity $\rdta(\xorg)$.

Note the $\eps'$ factor in the query complexity in Drucker's theorem.
When $\eps'$ is a constant close to $1/2$, Drucker's lower bound is
stronger than ours by a large constant factor.  However, when $\eps' = o(1)$,
his bound degrades significantly.  Couched in our notation, Drucker's
XOR Lemma yields
$\rdt_\eps(\xorg) = \Omega(\eps'k \rdt_{\eps'}(g))$, for some
$\eps' = O(\eps/k)$.  This simplifies to
$\rdt_\eps(\xorg) = \Omega(\eps R_{\eps/k}(g))$, a loss of a
factor of $k$.

As far as we know, it remains open whether this $\eps'$ factor is
needed in the query complexity lower bound of Drucker's XOR Lemma.
However, Shaltiel's counterexample~\cite{Shaltiel03} shows that the
$\eps'$ factor is required for distributional query complexity.  This
rules out the most direct approach for proving a tighter XOR Lemma for
$\rdt(\xorg)$.

Our %
paper %
is most closely related to that of Blais and
Brody~\cite{BlaisB19}, who give a strong direct sum theorem for
the expected query complexity of computing $k$ copies of $f$ in parallel, for any partial function $f$%
%
%
%
%
%
, and explicitly
conjecture that $\xor$ admits a strong direct sum theorem.
Both~\cite{BlaisB19} and our %
paper %
use techniques similar to work of
Molinaro et al.~\cite{MolinaroWY13, MolinaroWY15} who give strong
direct sum theorems for communication complexity.

Our strong direct sum theorem for $\xor$ is an example of a \emph{composition
  theorem}---a lower bound on the query complexity of functions of the
form $f\circ g$.  Several %
recent articles study %
composition theorems in query complexity.  Bassilakis et
al.~\cite{BassilakisDGHMT20} show that
$\rdt(f\circ g) = \Omega(\fbs(f)\rdt(g))$, where $\fbs(f)$ is the
\emph{fractional block sensitivity} of $f$.  Ben-David and
Blais~\cite{BenDavidB20a, BenDavidB20b} give a tight lower bound on
$\rdt(f\circ g)$ as a product of $\rdt(g)$ and a new measure they
define called \emph{$\noisyR(f)$}, which measures the complexity of
computing $f$ on noisy inputs.  They also characterize $\noisyR(f)$ in
terms of the gap-majority function.  Ben-David et
al~\cite{BenDavidGKW20} explicitly consider strong direct sum theorems
for composed functions in randomized query complexity, asking whether
the naive success amplification algorithm is necessary to compute
$f\circ g$.  They give a partial strong direct sum theorem, showing
that there exists a partial function $g$ such that computing $\xorg$
requires success amplification, even in a model where the abort
probability may be arbitrarily close to 1.\footnote{In this query
  complexity model, called $\PostBPP$, the query algorithm is allowed
  to abort with any probability strictly less than 1.  When it %
  does not %
  abort, it must output $f$ with probability at least $1-\eps$.}
Ben-David et al. explicitly ask whether there exists a total function
$g$ such that $\rdt(\xorg) = \Omega(k\log(k)\rdt(g))$.

%
\subsection{Our technique}    %
Our technique most closely follows the
strong direct sum theorem of Blais and Brody.  We start with a query
algorithm that computes $\xorg$ and use it to build a query algorithm
for computing $g$ with low error.  To do this, %
we will %
take an input for
$g$ and \emph{embed} it into an input for $\xorg$.  Given
$x \in \bitsn$, $i \in [k]$, and $y \in \bit^{n\times k}$, let
$\ygetix \deq (y^{(1)}, \ldots, y^{(i-1)}, x, y^{(i+1)}, \ldots
y^{(k)})$ denote the input obtained from $y$ by replacing the $i$-th
coordinate $y^{(i)}$ with $x$.  Note that if $x \sim \mu$ and
$y \sim \mu^k$,\footnote{We use $\mu^k$ to denote the distribution
  obtained on $k$-tuples of $\bitsn$ obtained by sampling each
  coordinate independently according to $\mu$.} then
$\ygetix\sim \mu^k$ for all $i \in [k]$.

We require the following observation %
\cite[Lemma 3.2]{Drucker12}.
\begin{lemma}[Drucker]
  Let $y\sim \mu^k$ be an input for a query algorithm $\cA$, and
  consider any execution of queries by $\cA$.  The distribution of
  coordinates of $y$, conditioned on the queries made by $\cA$,
  remains a product distribution.
\end{lemma}

In particular, the answers to $g(y^{(i)})$ remain independent bits
conditioned on any set of queries made by the query algorithm.
Our first observation is that in order to compute $\xorg(y)$ with high
probability, we must be able to compute $g(y^{(i)})$ with very high
probability for many $i$'s.  The intuition behind this observation is
captured by the following simple fact about the $\xor$ of independent
random bits.

Define the \emph{bias} of a random bit $X \in \bit$ as
$r(X) \deq \max_{b\in\bit}\Pr[X=b]$.  Define the \emph{advantage} of
$X$ as $\adv(X) \deq 2r(X)-1$.  Note that when $\adv(X) = \delta$,
then $r(X) = \frac{1}{2}(1+\delta)$.

\begin{fact}\label{fact:xor-adv}
  Let $X_1,\ldots, X_k$ be independent random bits, and let $a_i$ be
  the advantage of $X_i$.  Then, $$\adv(X_1\oplus \cdots \oplus X_k) =
  \prod_{i=1}^k \adv(X_i)\ .$$
\end{fact}
\begin{proof}
  For each $i$, let $b_i \deq \argmax_{b\in \bit} \Pr[X_i = b]$ and
  $\delta_i \deq \adv(X_i)$.  Then $\Pr[X_i = b_i] =
  \frac{1}{2}(1+\delta_i)$.
  We prove \expref{Fact}{fact:xor-adv} by induction on $k$.  When $k=1$,
  there is nothing to prove.  For $k=2$, note that
  \begin{align*}
    \Pr[X_1 \oplus X_2 = b_1\oplus b_2] &=
    \frac{1}{2}(1+\delta_1)\frac{1}{2}(1+\delta_2) +
    \frac{1}{2}(1-\delta_1)\frac{1}{2}(1-\delta_2) \\
    &= \frac{1}{4}(1+\delta_1+\delta_2 + \delta_1\delta_2) + \frac{1}{4}(1-\delta_1-\delta_2 + \delta_1\delta_2) \\
    &= \frac{1}{2}(1+\delta_1\delta_2)\ .
  \end{align*}
  Hence $X_1\oplus X_2$ has advantage $\delta_1\delta_2$ and the claim
  holds for $k=2$.
  For an induction hypothesis, suppose that the claim holds for
  $X_1\oplus \cdots \oplus X_{k-1}$.  Then, setting $Y \deq X_1 \oplus
  \cdots \oplus X_{k-1}$, by the induction hypothesis, we have $\adv(Y)
  = \prod_{i=1}^{k-1} \adv(X_i)$.  Moreover, $X_1\oplus \cdots \oplus
  X_k = Y\oplus X_k$, and
\begin{equation*}              %
 \adv(X_1 \oplus \cdots \oplus X_k) =
 \adv(Y\oplus X_k) = \adv(Y)\adv(X_k) = \prod_{i=1}^k \adv(X_i)\ . \qedhere
\end{equation*}
\end{proof}


Given an algorithm for $\xorg$ that has error $\eps$, it follows that
for typical leaves the advantage of computing $\xorg$ is
$\gtrsim 1-2\eps$.  \expref{Fact}{fact:xor-adv} shows that for such
leaves, the advantage of computing $g(y^{(i)})$ for most coordinates
$i$ is $\gtrsim (1-2\eps)^{1/k} = 1-\Theta(\eps/k)$.  Thus,
conditioned on reaching this leaf of the query algorithm, we could
compute $g(y^{(i)})$ with very high probability.  %
We would %
like to fix a
coordinate $i^*$ such that for most leaves, our advantage in computing
$g$ on coordinate $i^*$ is $1-O(\eps/k)$.  There are other
complications, namely that (i) our construction needs to handle aborts
gracefully and (ii) our construction must ensure that the algorithm
for $\xorg$ %
does not %
query the $i^*$-th coordinate too many times.  Our
construction identifies a coordinate $i^*$ and a string
$z \in \bit^{n\times k}$, and on input $x \in \bitsn$ it emulates a
query algorithm for $\xorg$ on input $\zgetsix$ and outputs our best
guess for $g(x)$ (which is now $g$ evaluated on coordinate $i^*$ of
$\zgetsix$), aborting when needed \eg, when the algorithm for
$\xorg$ aborts or when it queries too many bits of $x$.  We defer
full details of the proof to \expref{Section}{sec:main}.


%
\subsection{Preliminaries and notation}   %

A partial Boolean function on the domain $\bitsn$ is a function
$f: S \rightarrow \bit$ for some subset $S \subseteq \bitsn$.  Call
$S$ the set of valid inputs for $f$.  Let $f$ be a partial Boolean
function on $\bitsn$ and $\mu$ a distribution whose support is a
subset of the valid inputs. %
                            %
                            %
                            %
%
%
%
%
%
%
%
%
%
%
We use $[n]$ to denote the set
$\{1,\ldots, n\}$ and $X \in_R S$ to denote an element $X$ sampled uniformly
from a set $S$.  Let $\mu^k$ denote the distribution obtained on
$k$-tuples of $\bitsn$ obtained by sampling each coordinate
independently according to $\mu$.

An algorithm $\cA$ is a $[q,\delta,\eps,\mu]$-distributional query
algorithm for $f$ if $\cA$ is a deterministic algorithm with query
cost $q$ that computes $f$ with error probability at most $\eps$ and
abort probability at most $\delta$ when the input $x$ is drawn from
$\mu$.\footnote{Note: in the literature, the error probability is
  sometimes defined as being conditioned on not aborting 
(\eg,\cite{BenDavidGKW20}).  We define
  the error probabilty without conditioning to match
%
article \cite{BlaisB19}   %
most closely related to our work.}

Our main theorem is a direct sum result for $\xorg$ for %
expected %
randomized query complexity; however, \expref{Lemma}{lemma:main} uses
distributional query complexity with aborts.  To translate between the
two, we need two results from Blais and Brody~\cite{BlaisB19} that
connect the query complexities in the randomized, %
expected %
randomized, and distributional query models.

\begin{fact}[\cite{BlaisB19}, Proposition 14]
\label{fact:abort-average}
For every
partial %
function $f : \bitsn \to \bit$, every
$0 \le \epsilon < \frac12$ and every $0 < \delta < 1$,
\[
\delta \cdot \rdtde(f) \le \overline{\rdt}_\epsilon(f) \le \tfrac{1}{1-\delta}
\cdot \rdt_{\delta,(1-\delta)\epsilon}(f).
%
%
\]
\end{fact}


\expref{Fact}{fact:abort-average} shows that when $\delta = 1-\Omega(1)$, to
achieve a lower bound for $\rdta_\eps(f)$, it suffices to lower bound
$\rdt_{\delta,\eps}(f)$.
Next, we need the following generalization of Yao's minimax lemma,
which connects randomized and distributional query complexity in the
presence of aborts.
\begin{fact}[\cite{BlaisB19}, Lemma 15]
\label{fact:yao-abort}
For any $\alpha,\beta>0$ such that $\alpha+\beta \leq 1$, we have
\[
\max_\mu \ddt_{\delta/\alpha, \eps/\beta}^\mu(f) \leq \rdtde(f)
    \leq \max_{\mu} \ddt_{\alpha\delta, \beta\eps}^\mu(f).
\]
\end{fact}
For simplicity, it might be helpful to consider the simplest case
where $\alpha = \beta = \frac{1}{2}$.  In this case, we recover
$\max_\mu \ddt_{2\delta, 2\eps}^\mu(f) \leq \rdtde(f) \leq \max_{\mu}
\ddt_{\delta/2, \eps/2}^\mu(f)$.  \expref{Fact}{fact:yao-abort} shows that
to prove a lower bound on $\rdt_{\delta,\epsilon}(f)$, it suffices to
prove a lower bound on distributional complexity (albeit with a
constant factor increase in abort and error probabilities).

%
We will also use the following convenient facts about expected value.
\begin{fact}[Law of Conditional Expectations]\label{fact:prob-ex2}
  Let $X$ and $Y$ be random variables.  Then, we have $$\E[X] = \E[\E[X|Y]]\ .$$
\end{fact}

%
%
%
%
%
%
%
%
%
%
%
%
%
%
%
%
%
%
%
%

\begin{fact}[Markov Inequality for Bounded Variables]\label{fact:bounded-markov}
  Let $X$ be a real-valued random variable with $0\leq X \leq 1$.
  Suppose that $\E[X] \geq 1-\eps$.  Then, for any $T>1$ it holds that
  $$\Pr[X<1-T\eps] < \frac{1}{T}\ .$$
\end{fact}
\begin{proof}
  Let $Y \deq 1-X$.  Then, $\E[Y] \leq \eps$.  By Markov's Inequality
  we have
\begin{equation*}
  \Pr[X<1-T\eps] = \Pr[Y > T\eps] \leq \frac{1}{T}\ .\qedhere
\end{equation*}
\end{proof}

%
\section{Strong XOR lemma}\label{sec:main}   %

In this section, we prove our main result.

\begin{lemma}[Formal Restatement of \expref{Lemma}{lemma:main}]\label{lemma:main-technical}
  For every
  partial %
  function $g: \bitsn \rightarrow \bit$, every distribution
  $\mu$ on $\bitsn$, every $0 \leq \delta \leq \frac{1}{5}$, and every
  $0 < \eps \leq \frac{1}{200}$, we
  have
  %
  %
  $$\ddt_{\delta,\eps}^{\mu^k}(\xorg) \geq \frac{k}{25} \ddt_{\delta',\eps'}^\mu(g)\ ,$$ %
  $\delta' = 0.36+3\delta$ and $\eps' = \frac{15000\eps}{k}$.
\end{lemma}

\begin{proof}
  Let $q \deq \ddt_{\delta,\eps}^{\mu^k}(\xorg)$, and suppose that $\cA$
  is a $[q,\delta,\eps,\mu^k]$-distributional query algorithm for
  $\xorg$.  Our goal is to construct an
  $\left[O(q/k), \delta', \eps', \mu\right]$-distributional query
  algorithm for $g$.  Towards that end, for each leaf $\ell$ of
  $\cA$ define
  \begin{align*}
    b_\ell &\deq \argmax_{b \in \bit} \Pr_{x \sim \mu^k}[\xorg(x) = b | \leaf(\cA, x) = \ell] \\
    r_\ell &\deq \Pr_{x \sim \mu^k}[\xorg(x) = b_\ell | \leaf(\cA,x) = \ell] \\
    a_\ell &\deq 2r_\ell - 1\ .
  \end{align*}
  Call $a_\ell$ the \emph{advantage} of $\cA$ on leaf $\ell$.

  The purpose of $\cA$ is to compute $\xorg$; however, %
  we will %
  show that
  $\cA$ must additionally be able to compute $g$ reasonably
  well on many coordinates of $x$.  For any $i \in [k]$ and any leaf $\ell$, define
  \begin{align*}
    b_{i,\ell} &\deq \argmax_{b \in \bit} \Pr_{x \sim \mu^k}[b = g(x^{(i)}) | \leaf(\cA,x) = \ell] \\
    r_{i,\ell} &\deq \Pr_{x\sim \mu^k}[b_{i,\ell} = g(x^{(i)}) | \leaf(\cA,x) = \ell] \\
    a_{i,\ell} &\deq 2r_{i,\ell} - 1\ .
  \end{align*}

  If $\cA$ reaches leaf $\ell$ on input $y$, then write $\cA(y)_i \deq
  b_{i,\ell}$.  $\cA(y)_i$ represents $\cA$'s best guess for $g(y^{(i)})$.

  Next, we define some structural characteristics of leaves that %
  we will %
  need to complete the proof.

  \begin{definition}[Good leaves, good coordinates]
    \mbox{ }
    
    \begin{itemize}
    \item Call a leaf $\ell$ \emph{good} if $r_\ell \geq 1-50\eps$.
      Otherwise, call $\ell$ \emph{bad}.
    \item Call a leaf $\ell$ \emph{good for $i$} if
      $a_{i,\ell} \geq 1-5000\eps/k$.  Otherwise, call a leaf $\ell$
      \emph{bad for $i$}.
    \end{itemize}
  \end{definition}
  When a leaf is good for $i$, then $\cA$, conditioned on reaching
  this leaf, computes $g(x^{(i)})$ with very high probability.  Before
  presenting the main reduction, we give a few simple claims to help
  our proof.  Our first claim shows that we reach a good leaf with
  high probability.

  \begin{claim}\label{claim:leaf-good}
    $\Pr_{x \sim \mu^k}[\leaf(\cA,x) \text{ is bad } | \cA(x) \text{
      doesn't abort}] \leq \frac{1}{25}$.
  \end{claim}
  \begin{proof}
    Conditioned on $\cA$ not aborting, it outputs the correct value of
    $\xorg$ with probability at least
    $1-\frac{\eps}{1-\delta} \geq 1-2\eps$.  We analyze this error
    probability by conditioning on which leaf is reached.  Let $\nu$
    be the distribution on $\leaf(\cA,x)$ when $x \sim \mu^k$,
    conditioned on $\cA$ not aborting.  Let $L\sim \nu$.  Then, we
    have:
\begin{align*}
  1-2\eps &\leq \Pr_{x\sim \mu^k}[\cA(x) = \xorg(x) | \cA \text{ doesn't abort}] \\
  &= \sum_{\text{leaf }\ell } \Pr_{L \sim \nu}[L = \ell]\cdot \Pr[\cA(x) = \xorg(x) | L=\ell] \\
  &= \sum_\ell \Pr[L=\ell]\cdot r_\ell \\
  &= \E_L[r_L]\ .
\end{align*}

Thus, $\E[r_L] \geq 1-2\eps$.  Recalling that $\ell$ is good if
$r_\ell \geq 1-50\eps$ and using \expref{Fact}{fact:bounded-markov}, $L$
is bad with probability at most $\frac{1}{25}$.
  \end{proof}

  Next, we claim that each good leaf is good for many $i$.
  \begin{claim}\label{claim:leaf-good-i}
    Let $\ell$ be any good leaf, and let $I$ be uniform on $[k]$.
    Then, we have: $$\Pr_I[\ell \text{ is bad for } I] \leq \frac{1}{25}\ .$$
  \end{claim}
  \begin{proof}
    Fix a good leaf $\ell$, and let
    $\beta_\ell := \Pr_I[\ell \text{ is bad for } I]$.
    Recall that if $\ell$ is good, then $r_\ell \geq 1-50\eps$.
    Therefore, $a_\ell \geq 1-100\eps$.  Using $1+x\leq e^x$ and
    $e^{-2x} \leq 1-x$ (which holds for all $0\leq x \leq 1/2$), we
    have for any good leaf $\ell$
    $$1-100\eps \leq a_\ell = \prod_{i=1}^k a_{i,\ell} \\ \leq
    \left(1-\frac{5000\eps}{k}\right)^{k\beta_\ell} \\ \leq
    e^{-5000\eps\cdot \beta_\ell} \\
    \leq 1-2500\eps\beta_\ell\ .$$ Rearranging terms, we see that
    $\beta_\ell \leq \frac{1}{25}$.
  \end{proof}

  Next, we describe a randomized algorithm $\cA'$ for $g$ whose
  expected query cost, abort probability, and error probability match
  the guarantees we want to provide when the input $x \sim \mu$.
  %
  We will %
  complete the proof of \expref{Lemma}{lemma:main-technical} by
  fixing the randomness used in $\cA'$.  Our algorithm works by
  independently $z \sim \mu^k$ and  $i$ uniformly from $[k]$,
  embedding $x$ in the $i$-th coordinate of $z$, and emulating $\cA$
  on the resulting string.

  \begin{algorithm}
    \caption{$\cA'(x)$}
    \begin{algorithmic}[1]
      \State Independently sample $I$ uniformly from $[k]$ and  $z \sim \mu^k$.
      \State $y \gets \zgetsIx$
      \State Emulate algorithm $\cA$ on input $y$.
      \State Abort
      \begin{enumerate}[(i)]
      \item if $\cA$ aborts,
      \item if $\cA$ reaches a bad leaf, or
      \item if $\cA$ reaches a leaf that is bad for $I$.
      \item if $\cA$ queries more than
        $\frac{25q}{k}$ bits of $x$,
      \end{enumerate}
      \State
      Otherwise, output $\cA(y)$.
    \end{algorithmic}
  \end{algorithm}
  Note that the emulation is possible since whenever $\cA$ queries the
  $j$-th bit of $y^{(I)}$, we can query $x_j$, and we can emulate
  $\cA$ querying a bit of $y^{(i)}$ for $i\neq I$ directly since $z$
  is fixed.  We claim that (i) $\cA'$ makes at most $\frac{25q}{k}$
  queries, (ii) $\cA'$ aborts with probability at most
  $\delta + 0.12$, and (iii) $\cA'$ errs with probability at most
  $\frac{5000\eps}{k}$.

First, note that $\cA'$ makes at most $\frac{25q}{k}$ queries, since
it aborts instead of making more queries.

Second, consider the abort probability of $\cA'$.  Our algorithm
aborts if $\cA$ aborts, if we reach a bad leaf, if the leaf we reach
is bad for $I$, of if $\cA$ makes more than $\frac{25q}{k}$ bits of
$y^{(I)}$.  Let $\cE_1$ be the event that $\cA$ aborts on input $y$.
Similarly, let $\cE_2,\cE_3,\cE_4$ be the events that $\cA$ reaches a
bad leaf, $\cA$ reaches a leaf that is bad for $i$, and $\cA$ queries
more than $\frac{25q}{k}$ bits of $x$ respectively.  Since
$x\sim \mu$, $z \sim \mu^k$, and $I$ is uniform on $[k]$, it follows
that $y\sim \mu^k$.  By the abort guarantees of $\cA$, we have
$\Pr[\cE_1] \leq \delta$.  By \expref{Claim}{claim:leaf-good} we have
$\Pr[\cE_2 | \cE_1] \leq 1/25$, and by \expref{Claim}{claim:leaf-good-i}
we have $\Pr[\cE_3 | \cE_1,\cE_2] \leq 1/25$.  Thus, we have
$\Pr[\cE_1\vee \cE_2 \vee \cE_3] \leq \delta +\frac{2}{25}$.


Next, for each $i \in [k]$, let $q_i(y)$ denote the number of queries
that $\cA$ makes to $y^{(i)}$ on input $y$.  The query cost of $\cA$
guarantees that for each input $y$,
$\sum_{1\leq i\leq k} q_i(y) \leq q$.  Therefore, for any $y$, at most
$\frac{k}{25}$ indices $i\in [k]$ satisfy $q_i(y) \geq \frac{25q}{k}$.
Hence, for $I \in_R [k]$, $x\sim \mu$, and $z \sim \mu^k$, and
recalling that $y = \zgetsIx$, we have:
$\Pr[\cE_4] \leq \frac{1}{25}.$ By a union bound, we have
$\Pr_{I,z,x}[\cA' \text{ aborts on input } y] = \Pr[\cE_1 \vee \cE_2
\vee \cE_3 \vee \cE_4] \leq \delta + \frac{3}{25}=
\delta+0.12$.

Third, we analyze the error probability of $\cA'$.  This algorithm
errs only when it reaches a leaf that is good for $I$.  By
\expref{Claim}{claim:leaf-good-i}, we are correct with probability at
least
$r_{I,\ell} = \frac{1+a_{I,\ell}}{2} \geq 1-\frac{5000\eps}{k}.$
Thus, we have $\Pr[\cA' \text{ errs}] \leq \frac{5000\eps}{k}$.

%
Letting $X$ be the indicator variable for the event that $\cA'$ aborts
and $Y = (I,z)$, \expref{Fact}{fact:prob-ex2} gives
$$\Pr[\cA' \text{ aborts }] = \E[\cA' \text{ aborts }] =
\E_{I,z}[\E[\cA' \text{ aborts } | I,z]] = \E_{I,z}[\Pr[\cA' \text{
  aborts }]]\ .$$ Thus algortihm $\cA'$ is a randomized algorithm
that, when given an input $x \sim \mu$, makes at most $\frac{25q}{k}$
queries and has the following guarantees:

%
%
%
\begin{align*}
  \E_{I,z}[\Pr_x[\cA' \text{ aborts}]] = \Pr_{I,x,z}[\cA' \text{ aborts}] &\leq \delta + 0.12 \text{, and}\\
  \E_{I,z}[\Pr_x[\cA'(y)_{(I)} \neq g(x)]] = \Pr_{I,x,z}[\cA'(y)_{(I)} \neq g(x)] &\leq \frac{5000\eps}{k}\ .
\end{align*}
By Markov's inequality and a union bound, there must be a setting of
$(i^*,z^*)$ such that
$\Pr_x[\cA' \text{ aborts }] \leq 3\delta + 0.36$ and
$\Pr_x[\cA'(y)_{(i^*)} \neq g(x)] \leq \frac{15000\eps}{k}$.  Let
$\cA''$ be a \emph{deterministic} algorithm that takes an input
$x \sim \mu$ and emulates algorithm $\cA'$ with $i^*$ and $z^*$ in
place of the randomly sampled $I,z$. This algorithm queries at most
$\frac{25q}{k}$, aborts with probability at most $3 \delta + 0.36$,
and errs with probability at most $\frac{15000\eps}{k}$.  Thus, it is
a $[O(q/k), 3\delta+0.36, \frac{15000\eps}{k}, \mu]$-distributional
algorithm for $g$, as required.
\end{proof}

    
\subsection{Proof of \expref{Theorem}{thm:sds-xor}}
\begin{proof}[Proof of \expref{Theorem}{thm:sds-xor}.]
  Define $\eps' \deq 30000\eps$.  Let $\mu$ be the input distribution
  for $g$ achieving
  $\max_\mu \ddt_{\frac{1}{2}, \frac{\eps'}{k}}^\mu(g)$, and let $\mu^k$
  be the $k$-fold product distribution of $\mu$.
  By the first inequality of \expref{Fact}{fact:abort-average} and the
  first inequality of \expref{Fact}{fact:yao-abort}, we
  have
  $$\rdta_\eps(\xorg) \geq \frac{1}{50}\rdt_{\frac{1}{50},\eps}(\xorg) \geq \frac{1}{50}\ddt_{\frac{1}{25}, 2\eps}^{\mu^k}(\xorg)\ .$$
  Additionally, by \expref{Lemma}{lemma:main-technical} and the second
  inequalities of \expref{Fact}{fact:abort-average}
  and \expref{Fact}{fact:yao-abort}, we have
  $$\ddt_{\frac{1}{25}, 2\eps}^{\mu^k}(\xorg) \geq \frac{k}{120} \ddt_{\frac{1}{2},\frac{\eps'}{k}}^\mu(g) \geq \frac{k}{120}\rdt_{\frac{2}{3},\frac{4\eps'}{k}}(g) \geq \frac{k}{360}\rdta_{\frac{12\eps'}{k}}(g) \ .$$
  Thus, we have
  $\rdta_\eps(\xorg) =
  \Omega\left(\ddt_{\frac{1}{25},2\eps}^{\mu^k}(\xorg)\right)$ and
  $\ddt_{\frac{1}{25},2\eps}^{\mu^k}(\xorg) =
  \Omega\left(k\rdta_{\frac{12\eps'}{k}}(g)\right)$.  By standard
  success amplification
  $\rdta_{\frac{12\eps'}{k}}(g) = \Theta(\rdta_{\frac{\eps}{k}}(g))$.
  Putting these together yields
  $$\rdta_\eps(\xorg) =
  \Omega\left(\ddt_{\frac{1}{25},2\eps}^{\mu^k}(\xorg)\right) =
  \Omega\left(k\rdta_{\frac{12\eps'}{k}}(g)\right) =
  \Omega\left(\rdta_{\frac{\eps}{k}}(g)\right)\ ,$$ hence
  $\rdta_\eps(\xorg) = \Omega\left(k\rdta_{\frac{\eps}{k}}(g)\right)$ completing the proof.
\end{proof}



%% Theory of Computing
%% File: v019a011.bbl
%% Publication date: December 31, 2023
%% Authors: Joshua Brody, Jae Tak Kim, Peem Lerdputtipongporn, and Hariharan Srinivasulu
\newcommand{\prelim}{Preliminary version in\ }\newcommand{\prelims}{Preliminary
  versions in\ }
% Define empty bibhead if not already defined
\providecommand{\bibhead}[1]{}
% Define tocrefpdfbookmark if not already defined
\expandafter\ifx\csname pdfbookmark\endcsname\relax%
  \providecommand{\tocrefpdfbookmark}{}
\else\providecommand{\tocrefpdfbookmark}{%
   \hypertarget{tocreferences}{}%
   \pdfbookmark[1]{References}{tocreferences}}%
\fi

\tocrefpdfbookmark
\begin{thebibliography}{10}

\bibitem{BassilakisDGHMT20}\bibhead{BassilakisDGHMT20}
{\sc Andrew Bassilakis, Andrew Drucker, Mika G{\"{o}}{\"{o}}s, Lunjia Hu,
  Weiyun Ma, and Li-Yang Tan}: The power of many samples in query complexity.
\newblock In {\em Proc. 47th Internat. Colloq. on Automata, Languages, and
  Programming (ICALP'20)}, pp. 9:1--18. Schloss Dagstuhl--Leibniz-Zentrum fuer
  Informatik, 2020.
\newblock [\epfmtdoi{10.4230/LIPIcs.ICALP.2020.9}, \epfmt{arxiv}{2002.10654},
  \epfmt{eccc}{TR20-027}]

\bibitem{BenAsherN95}\bibhead{BenAsherN95}
{\sc Yosi Ben{-}Asher and Ilan Newman}: Decision trees with boolean threshold
  queries.
\newblock {\em J. Comput. System Sci.}, 51(3):495--502, 1995.
\newblock \prelim \href{https://doi.org/10.1109/SCT.1995.514729}{CCC'95}.
\newblock [\epfmtdoi{10.1006/jcss.1995.1085}]

\bibitem{BenDavidB20a}\bibhead{BenDavidB20a}
{\sc Shalev Ben{-}David and Eric Blais}: A new minimax theorem for randomized
  algorithms.
\newblock In {\em Proc. 61st FOCS}, pp. 403--411. IEEE Comp. Soc., 2020.
\newblock [\epfmtdoi{10.1109/FOCS46700.2020.00045}]

\bibitem{BenDavidB20b}\bibhead{BenDavidB20b}
{\sc Shalev Ben{-}David and Eric Blais}: A tight composition theorem for the
  randomized query complexity of partial functions.
\newblock In {\em Proc. 61st FOCS}, pp. 240--246. IEEE Comp. Soc., 2020.
\newblock [\epfmtdoi{10.1109/FOCS46700.2020.00031}, \epfmt{arxiv}{2002.10809}]

\bibitem{BenDavidGKW20}\bibhead{BenDavidGKW20}
{\sc Shalev Ben{-}David, Mika G{\"{o}}{\"{o}}s, Robin Kothari, and Thomas
  Watson}: When is amplification necessary for composition in randomized query
  complexity?
\newblock In {\em Proc. 24th Internat. Conf. on Randomization and Computation
  (RANDOM'20)}, pp. 28:1--16. Schloss Dagstuhl--Leibniz-Zentrum fuer
  Informatik, 2020.
\newblock [\epfmtdoi{10.4230/LIPICS.APPROX/RANDOM.2020.28},
  \epfmt{arxiv}{2006.10957}]

\bibitem{BenDavidK18}\bibhead{BenDavidK18}
{\sc Shalev Ben{-}David and Robin Kothari}: Randomized query complexity of
  sabotaged and composed functions.
\newblock {\em Theory of Computing}, 14(5):1--27, 2018.
\newblock \prelim
  \href{https://doi.org/10.4230/LIPIcs.ICALP.2016.60}{ICALP'16}.
\newblock [\epfmtdoi{10.4086/toc.2018.v014a005}, \epfmt{arxiv}{1605.09071},
  \epfmt{eccc}{TR16-087}]

\bibitem{BlaisB19}\bibhead{BlaisB19}
{\sc Eric Blais and Joshua Brody}: Optimal separation and strong direct sum for
  randomized query complexity.
\newblock In {\em Proc. 34th Comput. Complexity Conf. (CCC'19)}, pp. 29:1--17.
  Schloss Dagstuhl--Leibniz-Zentrum fuer Informatik, 2019.
\newblock [\epfmtdoi{10.4230/LIPIcs.CCC.2019.29}, \epfmt{arxiv}{1908.01020}]

\bibitem{Drucker12}\bibhead{Drucker12}
{\sc Andrew Drucker}: Improved direct product theorems for randomized query
  complexity.
\newblock {\em Comput. Complexity}, 21(2):197--244, 2012.
\newblock \prelim \href{https://doi.org/10.1109/CCC.2011.29}{CCC'11}.
\newblock [\epfmtdoi{10.1007/s00037-012-0043-7}, \epfmt{arxiv}{1005.0644},
  \epfmt{eccc}{TR10-080}]

\bibitem{ImpagliazzoRW94}\bibhead{ImpagliazzoRW94}
{\sc Russell Impagliazzo, Ran Raz, and Avi Wigderson}: A direct product
  theorem.
\newblock In {\em Proc. 9th IEEE Conf. Structure in Complexity Theory
  (SCT'94)}, pp. 88--96. IEEE Comp. Soc., 1994.
\newblock [\epfmtdoi{10.1109/SCT.1994.315814}]

\bibitem{JainKS10}\bibhead{JainKS10}
{\sc Rahul Jain, Hartmut Klauck, and Miklos Santha}: Optimal direct sum results
  for deterministic and randomized decision tree complexity.
\newblock {\em Inform. Process. Lett.}, 110(20):893--897, 2010.
\newblock [\epfmtdoi{10.1016/j.ipl.2010.07.020}]

\bibitem{MolinaroWY13}\bibhead{MolinaroWY13}
{\sc Marco Molinaro, David~P. Woodruff, and Grigory Yaroslavtsev}: Beating the
  direct sum theorem in communication complexity with implications for
  sketching.
\newblock In {\em Proc. 24th Ann. ACM--SIAM Symp. on Discrete Algorithms
  (SODA'13)}, pp. 1738--1756. SIAM, 2013.
\newblock [\epfmtdoi{10.1137/1.9781611973105.125}]

\bibitem{MolinaroWY15}\bibhead{MolinaroWY15}
{\sc Marco Molinaro, David~P. Woodruff, and Grigory Yaroslavtsev}:
  Amplification of one-way information complexity via codes and noise
  sensitivity.
\newblock In {\em Proc. 42nd Internat. Colloq. on Automata, Languages, and
  Programming (ICALP'15)}, pp. 960--972. Springer, 2015.
\newblock [\epfmtdoi{10.1007/978-3-662-47672-7\_78}, \epfmt{eccc}{TR15-031}]

\bibitem{NisanRS99}\bibhead{NisanRS99}
{\sc Noam Nisan, Steven Rudich, and Michael~E. Saks}: Products and help bits in
  decision trees.
\newblock {\em SIAM J. Comput.}, 28(3):1035--1050, 1998.
\newblock \prelim \href{https://doi.org/10.1109/SFCS.1994.365683}{FOCS'94}.
\newblock [\epfmtdoi{10.1137/S0097539795282444}]

\bibitem{Shaltiel03}\bibhead{Shaltiel03}
{\sc Ronen Shaltiel}: Towards proving strong direct product theorems.
\newblock {\em Comput. Complexity}, 12(1):1--22, 2003.
\newblock \prelim \href{https://doi.org/10.1109/CCC.2001.933878}{CCC'01}.
\newblock [\epfmtdoi{10.1007/s00037-003-0175-x}, \epfmt{eccc}{TR01-009}]

\end{thebibliography}

\end{document}
